% ---------------------------------------------------------------------------
% erw-cdr-slides.tex
% Two-slide summary of the ERW CDR modelling chain and its references.
% Build:  pdflatex erw-cdr-slides.tex
% ---------------------------------------------------------------------------
\documentclass[aspectratio=169,10pt]{beamer}

\usetheme{default}
\setbeamertemplate{navigation symbols}{}
\setbeamertemplate{footline}[frame number]
\setbeamerfont{frametitle}{size=\large,series=\bfseries}

\usepackage{amsmath}
\usepackage{booktabs}
\usepackage{xcolor}
\definecolor{erwblue}{HTML}{1F4E5F}
\definecolor{erwgreen}{HTML}{2C7A3F}
\setbeamercolor{frametitle}{fg=erwblue}
\setbeamercolor{structure}{fg=erwblue}

\newcommand{\co}{\mathrm{CO_2}}

\title{Enhanced Rock Weathering: CDR Modelling}
\date{}

\begin{document}

% ===========================================================================
% SLIDE 1 -- methodology / equations
% ===========================================================================
\begin{frame}{How the CDR model works}

\small
For every $\sim$1\,km$^2$ pixel the model asks: \emph{how much $\co$ does one tonne of
crushed basalt durably remove here?} A chemical ceiling is reduced by fractional
``haircuts,'' scaled by the application rate, then netted against the project's own
emissions.

\medskip
\textbf{1. Chemical ceiling (stoichiometry).} CaO/MgO each consume 2\,mol $\co$:
\[
\mathrm{CDR_{max}} =
\Big(\underbrace{0.10}_{\mathrm{CaO}}\cdot\tfrac{2\cdot44.01}{56.08}
   + \underbrace{0.07}_{\mathrm{MgO}}\cdot\tfrac{2\cdot44.01}{40.30}\Big)\cdot 1000
\;\approx\; 310~\mathrm{kg~\co/t~basalt}.
\]

\textbf{2--4. Fractional haircuts} (reactivity, grain size, climate, soil pH, carbonate re-precipitation):
\[
\mathrm{CDR_{per\,t}}(px)=\mathrm{CDR_{max}}\times
\underbrace{0.20}_{\text{reactive}}\times
\underbrace{\big(\tfrac{100}{d}\big)^{0.5}}_{\text{grain}}\times
\underbrace{f_{\mathrm{MAT}}f_{\mathrm{MAP}}f_{\mathrm{pH}}}_{\text{climate, capped}}\times
\underbrace{(1-p_{\mathrm{ped}})}_{\text{exported}} .
\]
\[
f_{\mathrm{MAT}}=\mathrm{clamp}\!\big(\tfrac{\mathrm{MAT}}{11^\circ\mathrm C},0.3,3\big),\quad
f_{\mathrm{MAP}}=\mathrm{clamp}\!\big(\tfrac{\mathrm{MAP}}{1000\,\mathrm{mm}},0.3,3\big),\quad
f_{\mathrm{pH}}:\text{triangular, peak}\approx\mathrm{pH}\,5.
\]

\textbf{5--7. Yield, net of emissions, revenue.}
\[
\mathrm{CDR_{gross}}=\mathrm{rate}\cdot\tfrac{\mathrm{CDR_{per\,t}}}{1000},\qquad
\mathrm{CDR_{net}}=\max\!\big(0,\;\mathrm{CDR_{gross}}-\mathrm{rate}\cdot\tfrac{\mathrm{LCA_{per\,t}}}{1000}\big),
\]
\[
\mathrm{revenue}=\mathrm{CDR_{net}}\cdot(\,p_{\mathrm{carbon}}-\mathrm{MRV}\,),\qquad
\text{NPV regime: $\times$ 5-yr phasing factor }\approx 0.79 .
\]

\medskip
\footnotesize\textcolor{erwgreen}{In short: chemistry sets the ceiling ($\sim$310\,kg/t);
geography (temperature, rainfall, acidity, aridity) decides how much of it you actually
capture ($\sim$80\,kg/t at the reference climate).}

\end{frame}

% ===========================================================================
% SLIDE 2 -- references
% ===========================================================================
\begin{frame}{Key references for methodology \& coefficients}
\scriptsize

\begin{columns}[T]
\begin{column}{0.5\textwidth}
\textbf{Chemistry \& stoichiometry}\\
Lewis et al.\ 2021, \emph{Appl.\ Geochem.}\ 132:105023 \\
Beerling et al.\ 2020, \emph{Nature} 583:242 (SI)

\medskip
\textbf{Reactive fraction (0.20)}\\
Lewis et al.\ 2021 --- load-bearing calibration

\medskip
\textbf{Grain size \& grinding ($\beta{=}0.5$, $\alpha{=}1.2$)}\\
Strefler et al.\ 2018, \emph{ERL} 13:034010

\medskip
\textbf{Climate scaling}\\
Kanzaki et al.\ 2023, \emph{PNAS Nexus} 2:pgad059 \\
Baek et al.\ 2023, \emph{Earth's Future} 11

\medskip
\textbf{pH dependence}\\
White \& Brantley 2003, \emph{Chem.\ Geol.}\ 202:479 \\
Brantley et al.\ 2008, \emph{Kinetics of Water-Rock Interaction}
\end{column}

\begin{column}{0.5\textwidth}
\textbf{Pedogenic-carbonate deduction}\\
Renforth 2019, \emph{Nat.\ Commun.}\ 10:1401 \\
Reershemius \& Suhrhoff 2024, \emph{GCB} 30:e17025 \\
UNEP 1992 (aridity bands); Zomer et al.\ 2022, \emph{Sci.\ Data} 9:409

\medskip
\textbf{Net CDR / lifecycle emissions}\\
Eufrasio et al.\ 2022, \emph{Commun.\ Earth Environ.}\ 3:106 \\
EEA EMEP 2023; UK BEIS 2023; Nemecek \& Schnetzer 2011

\medskip
\textbf{Time phasing \& MRV}\\
Kanzaki et al.\ 2023; Levy et al.\ 2024, \emph{ES\&T} 58:17215

\medskip
\textbf{Application rate \& targeting}\\
Beerling 2020; Baek 2023; Kantola et al.\ 2023, \emph{GCB} 29:7012 \\
Aramburu Merlos et al.\ 2023, \emph{Geoderma} 432:116421

\vspace{1em}
\textcolor{erwgreen}{Full catalogue (S1--S37) with DOIs: \texttt{docs/erw-parameters.md}}
\end{column}
\end{columns}

\end{frame}

\end{document}
